\section{Resultados}
\subsection{Fabricación del primer prototipo}
Para el primer prototipo se tomaron en cuenta requerimientos
sujetos a nuestras posibilidades de fabricación y recursos destinados tras
haber seleccionado la instrumentación, el microcontrolador y demas partes motrices.
\\
En este orden de ideas se empleó la fabricación de las partes con diseño propio,
mediante impresión 3D, la base del prototipo se creó empleando una tabla de MDF reciclado
de la carpintería de la universidad donde además se le hicieron los agujeros respectivos. 
Adicionalmene para los brazos AB y BC se cortó una varilla de aluminio a la medida deseada.
\\ Para el sistema motriz se emplearon partes comunes en máquinas CNC domésticas (como impresoras 3D),
además de rodamientos comunes (del tipo 625 y 608)  en los sitios requeridos.
\\
\\
Todas estas características resultan en un sistema altamente fiable, rápido de armar y fácil de controlar;
El cual adicionalmente minimiza los factores de incertidumbre en el movimiento (tales como la fricción y la falta de rígidez de las partes), 
permitiendo así mediciones precisas.

\subsection{Implementación de la instrumentación}

